\section{Conclusions and Next Steps}
\label{sec:conclusions}

The main structural conclusion of the project is that the experiments should no
longer be presented as a single linear sequence of QAE trials. They form three
interlocking workflows.

\begin{itemize}[leftmargin=*]
  \item \textbf{Workflow A} studies angle-structured QAE integration of a
        function \(g\) under uniform quadrature. This is where the original
        \(g_0,g_1,g_2\) hierarchy, MLAE/IQAE/canonical-QAE comparisons, and the
        forthcoming quadrature-rule experiments belong.
  \item \textbf{Workflow B} studies Grover--Rudolph preparation of a finite
        probability law \(p\) and direct computation of finite-law quantities
        \(I_p(f)\). This includes the distribution-validation campaigns,
        affine and beta-bump scaling, MC/QMC baselines, the \(n=6\) affine
        result, and the centered Gaussian test.
  \item \textbf{Workflow C} is the backend-aware control layer. It explains the
        role of local audits, compilation choices, noise-model simulations,
        mitigation, layout selection, repetition rules, and rejection of
        amplified circuits whose predicted or measured behaviour is not stable.
\end{itemize}

This organization also resolves a conceptual ambiguity in the earlier draft.
The calibration functions \(g_0,g_1,g_2\) are used as integrands in Workflow A.
The later Grover--Rudolph profiles such as \texttt{gr\_affine\_n4},
\texttt{gr\_affine\_n6}, and
\texttt{gr\_gaussian\_mid\_sigma025\_n6} are used to define probability laws in
Workflow B; the quantities of interest are then classical test functions
integrated against those laws.

Phase 1 establishes a complete and reproducible first benchmark cycle on IBM
hardware. The main conclusions are:
\begin{itemize}[leftmargin=*]
  \item local audit before QPU execution is essential; it prevented blind
        escalation to deeper circuits and made the \(g_0,k=2\) anomaly
        interpretable;
  \item compilation settings materially affect hardware accuracy, even when the
        logical and saved transpiled circuits are ideal-correct;
  \item optimization level 2 provided the accepted evidence chain on
        \texttt{ibm\_fez};
  \item after compilation tuning, the affine and quadratic benchmarks both pass
        through \(k=2\) and produce accurate MLAE estimates.
\end{itemize}

The accepted Phase 1 result set should be treated as the baseline for the next
stage. Further QPU spending should pause until Phase 2 is designed explicitly.
The most natural Phase 2 directions are:
\begin{enumerate}[leftmargin=*]
  \item compare layout-selection strategies for the same \(g_0,g_1,g_2\)
        schedules;
  \item add readout mitigation and zero-noise extrapolation to quantify whether
        the residual \(g_0,k=2\) bias can be reduced;
  \item extend the benchmark from the SISC calibration functions to larger
        structured Grover--Rudolph instances;
  \item convert the current report into an article-ready methodology paper,
        emphasizing anomaly detection, compilation-sensitive correction, and
        reproducible hardware evidence.
\end{enumerate}

The immediate Phase 2 action is therefore a credit-free layout and compilation
audit over optimization levels and transpiler seeds on \texttt{ibm\_fez}. The
audit ranks candidates by backend-noise deviation, two-qubit-gate count, and
depth before any additional QPU jobs are considered.

The first Phase 2 audit identified an actionable \(g_0,k=2\) candidate with
optimization level 2 and transpiler seed 67890. The candidate has 19 two-qubit
gates, depth 76, and backend-noise predicted deviation 0.008545, compared with
the accepted Phase 1 hardware deviation 0.047363. This motivates a two-circuit
microcampaign repeating only \(g_0,k=0\) and \(g_0,k=2\). The strong \(g_1\) and
\(g_2\) Phase 1 results should not be repeated unless mitigation or larger
instances are introduced.

The Phase 2 microcampaign did not confirm the audit prediction. The measured
\(g_0,k=2\) probability was 0.306641, with absolute deviation 0.056641, compared
with the audit prediction 0.258545 and deviation 0.008545. The corresponding
MLAE estimate was \(\hat{a}=0.239373\), with absolute error 0.010627. This is
slightly worse than the accepted Phase 1 \(g_0\) repeat. The conclusion is that
the backend noise model used in the audit is not sufficiently predictive for
this sensitive circuit. Further QPU spending on random seed or layout searches
is therefore not justified; the next Phase 2 step should be explicit mitigation.
The next prepared microcampaign uses Qiskit Runtime dynamical decoupling and
Pauli/measurement twirling on only the \(g_0\) calibration pair, keeping the
accepted Phase 1 compilation seed fixed so that mitigation is the experimental
variable.

The mitigation microcampaign improved the \(g_0\) calibration result. With
dynamical decoupling and gate/measurement twirling, \(g_0,k=2\) gave
\(\hat{p}=0.291016\), absolute deviation 0.041016, and MLAE estimate
\(\hat{a}=0.242160\) with absolute error 0.007840. This is the best \(g_0\)
result obtained so far, improving on the accepted Phase 1 repeat
\(0.008949\) and the Phase 2 layout-seed repeat \(0.010627\). The circuit
nevertheless remains in the watch region, so the project should treat \(g_0\) as
a sensitive calibration case and use the robust \(g_1/g_2\) results as the main
positive hardware evidence.

Phase 1b simulator experiments support MLAE as the main hardware workflow. The
IQAE-style adaptive proxy reduces ideal statistical error, but at a substantially
larger oracle-query budget. Canonical QAE and MC/QMC are degenerate on the
current four-point calibration family, so the next fair comparison should use
non-special amplitudes or larger structured Grover--Rudolph instances.

The non-degenerate simulator extensions support the same decision. Across
synthetic amplitudes and larger midpoint-grid functions, the IQAE proxy improves
ideal MAE by roughly a factor of two, but typically requires about twice the
oracle-query budget of MLAE \(K=\{0,1,2\}\) and introduces adaptive execution
overhead. MLAE therefore remains the preferred near-term hardware method.

The project has now reached the next methodological checkpoint: larger
Grover--Rudolph state preparation has been validated directly before QAE is
attached to it. The Phase 0-GR audit shows that the paper-aligned
uniformly-controlled-\(R_y\) construction is viable on \texttt{ibm\_fez} for
eight- and sixteen-point dyadic grids. The first hardware distribution campaign
then measured \texttt{gr\_sin2\_n3}, \texttt{gr\_sin2\_half\_n3}, and
\texttt{gr\_affine\_n4}. The two eight-point circuits achieved TVD 0.0313 and
0.0216, while the sixteen-point affine circuit achieved TVD 0.0479. These data
support the state-preparation workflow, but also show that the first
sixteen-point case is the right object for a focused repeat. That repeat,
performed with 4096 shots, dynamical decoupling, and twirling, reduced TVD to
0.0341 and the maximum per-cell deviation to 0.0086. This is sufficient evidence
to prepare the first amplified test, but it should remain minimal:
\texttt{gr\_affine\_n4}, marked event \(x\ge 1/2\), and MLAE schedule
\(K=\{0,1\}\). The first amplified readout confirms that this caution was
necessary: \(k=0\) remains compatible with the distribution-measurement scale,
but \(k=1\) gives \(\hat p_1=0.173828\) instead of the expected 0.052765. The
resulting MLAE error is 0.059854. The \(k=2\) amplified circuit should therefore
remain deferred. The next action is a credit-free diagnosis of
\texttt{gr\_affine\_n4}, \(k=1\), searching for a compilation/layout candidate
whose ideal and noisy predictions are consistent before any repeat is submitted.
That diagnosis identified an opt-3, seed-12345 candidate with depth 152 and 35
two-qubit gates, compared with depth 225 and 51 two-qubit gates in the executed
opt-2 circuit. A single-circuit repeat of \(k=1\) with this candidate is
justified; \(k=2\) remains deferred. The repeat improved the hardware probability
from 0.173828 to 0.124023, but the remaining deviation, 0.071259, is still too
large. The next step should use the existing counts to test readout
postprocessing and event-level bias before any further QPU submission. That
test has now been performed with the real transpiled event-qubit layouts. The
event bit is measured on \(q[137]\) in the opt-2 circuit and on \(q[134]\) in
the opt-3 repeat. Current backend assignment probabilities are below the level
needed to explain the observed lower-to-upper leakage; readout correction leaves
the opt-3 marked probability essentially unchanged, 0.124023 to 0.124199.
Therefore the next investigation should target circuit-level leakage,
amplification sensitivity, and mitigation beyond final readout assignment,
while \(k=2\) remains deferred. A subsequent block-level noisy simulation
confirms this direction: the Grover--Rudolph loading block \(A\) remains stable,
but the full \(AS_0A^\dagger OA\) iterate moves the event mass from the ideal
0.052765 to 0.112793 under the backend-noise model. The practical next step is
therefore to study lower-depth or partially mitigated amplification blocks
before any additional hardware campaign. The first variant audit has already
ruled out simple replacements of the zero reflection: \texttt{mcp}, diagonal,
and direct-tree variants are all worse than the current \texttt{ucry}+\texttt{mcx}
opt-3 candidate. Consequently, further progress likely requires a different
amplification strategy or a smaller event/circuit family, not just a different
decomposition of the same \(k=1\) Grover iterate.

The constructive alternative is to use the measured Grover--Rudolph
distribution directly. For the same \texttt{gr\_affine\_n4} upper-half event,
the mitigated distribution estimate gives 0.671387 against the exact value
0.680380, with absolute error 0.008993 and a Wilson interval containing the
truth. This is substantially better than the available amplified MLAE estimate.
The immediate project decision is therefore to treat direct mitigated
Grover--Rudolph distribution estimation as the near-term hardware workflow for
sixteen-point grids, while amplified MLAE remains a diagnostic target requiring
additional circuit-level stabilization. The functional-level analysis refines
this decision: event probabilities and smooth test-function expectations are already useful,
whereas right-tail-sensitive moments and payoff-like functions still reveal
residual distribution bias. Future hardware campaigns should therefore report
quantity-specific error, not only global TVD or QAE amplitude error.
Classical postprocessing can quantify this bias, but the first leave-one-out
test shows that a single affine correction across test functions is not robust.
The next calibration layer should therefore be functional-aware or
distribution-structure-aware, rather than a global scalar correction.

The Phase 5 mitigated distribution suite strengthens this decision while adding
one caution. The deferred localized profile \texttt{gr\_beta\_bump\_n4} is
hardware-viable as a direct distribution experiment, with TVD 0.0456 and maximum
cell deviation 0.0241, but it is not yet clean enough to attach amplified QAE.
The two eight-point repeats show that Runtime mitigation is not uniformly
beneficial: it is neutral for \texttt{gr\_sin2\_half\_n3} and slightly worsens
\texttt{gr\_sin2\_n3}. A follow-up audit and focused repeat of the exact
\texttt{gr\_beta\_bump\_n4} law, \(F(x)=x(1-x)^4\), further supports caution.
Although the audit identified an opt-3, seed-45678 candidate with predicted TVD
0.0217, the hardware repeat gave TVD 0.0502. The maximum cell deviation improved
from 0.0241 to 0.0170, but errors for quantities of interest were mixed and most did not
improve. The next experimental step should therefore not be a deeper Grover
schedule. Amplified QAE remains deferred for localized profiles; the productive
path is direct distribution diagnostics, quantity-specific error reporting,
and model-aware postprocessing.

Phase 6 begins that postprocessing program. The first extended distribution
analysis adds classical fidelity, Hellinger distance, Euclidean distance, and
bootstrap intervals to TVD. These metrics reinforce the current decision:
\texttt{gr\_affine\_n4} benefits from mitigation, \texttt{gr\_sin2\_n3} does
not, and the \texttt{gr\_beta\_bump\_n4} repeat is not a robust improvement even
though its largest cell error is smaller. The simulator-only purification
diagnostic inspired by mixed-state preparation confirms that the
Grover--Rudolph loader can be viewed as an eigenvalue-encoding layer with an
optional entropy-injection register. This is a useful conceptual extension, but
not yet a reason to spend QPU credits. The Phase 6 decision table therefore
promotes only direct distribution and selected quantity-of-interest estimates:
\texttt{gr\_sin2\_half\_n3} is a direct distribution result,
\texttt{gr\_sin2\_n3} is a direct result with caveats, and mitigated
\texttt{gr\_affine\_n4} supports selected functionals. Localized
\texttt{gr\_beta\_bump\_n4} profiles remain diagnostic only. No
Grover--Rudolph profile is currently promoted to amplified QAE on
\texttt{ibm\_fez}. At the functional level, however, the data already contain a
set of defensible direct estimates: five case--functional pairs have absolute
error below \(5\times 10^{-3}\) with bootstrap coverage of the exact value, and
seven more have error below \(10^{-2}\) with coverage. This gives a clear
near-term publication path based on direct Grover--Rudolph probability loading,
metric-rich distribution benchmarking, and quantity-of-interest postprocessing.

The Phase 7 higher-shot repeats close the first hardware campaign. They show
that more shots alone do not turn every measured distribution into a stable
global probability-preparation claim. Among the final direct-distribution
results, only mitigated \texttt{gr\_affine\_n4} is promoted as a
paper-quality global distribution result, with TVD 0.028017 and maximum cell
deviation 0.008197. The localized \texttt{gr\_beta\_bump\_n4} profile remains a
backend diagnostic distribution, with TVD 0.057598, despite having two useful
functionals below \(10^{-2}\). The two \(n=3\) raw repeats also reinforce the
same distinction: several algebraic functionals remain useful, but the global
probability laws are not uniformly stable under repetition.

The final paper-quality set should therefore be stated at two levels. At the
distribution level, the defensible hardware claim is the mitigated
\texttt{gr\_affine\_n4} loader. At the algebraic-probability level, the strongest
functionals are \(I_p(\sin(\pi x))\) for \texttt{gr\_affine\_n4} and \(I_p(x)\)
for \texttt{gr\_sin2\_n3}. A second tier of useful results, including selected
call-payoff and moment functionals, should be reported with explicit bias
caveats. All rejected or diagnostic entries remain scientifically important:
they identify backend drift, localized-profile sensitivity, and event-tail
bias, which are precisely the effects that would be hidden by reporting only a
single QAE amplitude error.

The first cross-backend Phase 8 test strengthens this conclusion. Repeating the
mitigated \texttt{gr\_affine\_n4} direct-distribution campaign on
\texttt{ibm\_kingston} improves the global distribution metrics relative to
\texttt{ibm\_fez}: TVD decreases from 0.028017 to 0.021823, Hellinger distance
from 0.023874 to 0.019060, and classical fidelity increases from 0.998860 to
0.999274. Kingston also reduces the absolute error for every selected
functional, including the previously diagnostic upper-half event. The upper-half
event is still not promoted to the primary paper-quality set, but the systematic
improvement confirms that a nontrivial part of the residual error is backend
dependent rather than intrinsic to the Grover--Rudolph circuit.

The campaign should now proceed selectively. The affine result is strong enough
to serve as the cross-backend anchor for the paper. Further QPU spending should
focus on one of two sharply defined questions: whether the \(n=3\) drift seen on
\texttt{ibm\_fez} also appears on \texttt{ibm\_kingston}, or whether the
localized \texttt{gr\_beta\_bump\_n4} instability is backend-specific. Amplified
QAE should remain deferred until these direct-distribution cross-backend
controls are understood.

The next methodological step is Phase 9: integration scaling. The goal is to
test whether the direct Grover--Rudolph workflow remains useful when the
finite probability law is refined from \(2^3\) and \(2^4\) grid cells to
\(2^5\) cells. This phase explicitly separates discretization error,
state-loading error, finite-shot error, and backend-noise-induced error for
each quantity of interest. This separation is important for the paper: a
hardware result should not be interpreted as a quantum integration advantage
unless the classical grid error, the probability-loading error, and the
quantity-specific hardware bias are all reported. Phase 9 should therefore
start as a credit-free audit on \texttt{ibm\_kingston}, using the affine
cross-backend result as the anchor and treating \(n=5\) cases as candidates only
if their transpiled depth, two-qubit count, noisy TVD, and maximum
quantity-of-interest error remain below the predeclared thresholds.

The first Phase 9 hardware submission confirms that this route is viable. On
\texttt{ibm\_kingston}, the mitigated 32-cell affine law achieves TVD 0.033694,
Hellinger distance 0.031298, classical fidelity 0.998042, and maximum cell
deviation 0.008438. This promotes \texttt{gr\_affine\_n5} to the first
paper-quality scaling result at the distribution level. At the functional
level, however, only
\(I_p(\sin(\pi x))\) is cleanly compatible with the exact finite-law value, with
absolute error 0.003601 and bootstrap interval coverage. The call payoff is
usable with a bias caveat, while the mean, second moment, and upper-half event
remain diagnostic.

The localized 32-cell stress test then establishes the current boundary of the
method. Although the \texttt{gr\_beta\_bump\_n5} circuit compiles with depth 162
and 44 two-qubit gates, the hardware result gives TVD 0.073298, Hellinger
distance 0.110712, classical fidelity 0.975636, and no quantity-of-interest
interval coverage. This does not pass the Phase 9 distribution threshold. The
conclusion is therefore sharper than a uniform success or failure: direct
Grover--Rudolph sampling scales to 32 cells for the affine law, but localized
probability profiles remain a backend-sensitive stress regime. Amplified QAE
should still remain deferred; the next methodological step should compare this
QPU sampling workflow against classical MC/QMC baselines under the same finite
laws.

Phase 10 performs that comparison without spending additional QPU credits. It
shows that the 32-cell affine result is not uniformly competitive with MC at
the same sample count. The strongest finite-law integral remains
\(I_p(\sin(\pi x))\), where the QPU absolute error 0.003601 is comparable with
the MC RMSE 0.003294 for \(M=8192\). The mean, second moment, call payoff, and
upper-half mass are instead dominated by hardware bias relative to MC. The
localized beta-like law is worse than MC for every tested quantity and should
remain a diagnostic stress case. The QMC baseline, implemented through the exact
finite inverse CDF, is much more accurate than both QPU and MC once the law is
known. This reinforces the right interpretation of the project at this stage:
we have demonstrated a reproducible hardware workflow for preparing and
sampling structured finite laws, but not a quantum integration advantage over
classical sampling baselines.

The angle-structure classification should now be treated as a design variable
for the next campaign. The experiments support the distinction between
QAE-efficient integration and direct Grover--Rudolph distribution integration.
For QAE-efficient integration, the relevant object is the angle map
\(\Theta_g=2\arcsin\sqrt{g}\), and low-degree classes
\(\mathcal{G}^{(d)}_n\) are the natural first candidates because they correspond
to simpler amplitude-oracle structure. For direct distribution integration,
high angle degree of the test function \(f\) does not prevent estimating
\(I_p(f)\) from measured frequencies, but it is a warning that the same
quantity may be expensive or unstable if recast as an amplified QAE amplitude.
The next local phase should therefore classify every proposed probability
profile and every proposed test function before backend submission, then choose
between direct sampling and amplified QAE according to that structural
classification and the MC/QMC baseline scale.

This policy has now been made explicit in Phase 11. The decision table promotes
the 32-cell affine law with \(f(x)=\sin(\pi x)\) as the current paper-quality
direct-sampling result. It also identifies low-angle \(\sin^2(\pi x)\) pairs,
such as the upper-half event, as simulator-level amplified-QAE candidates, and
defers QPU repeats for localized beta-like tail quantities. This is the
strongest current use of the angle-structure theory in the experimental
workflow: it does not claim that low angle degree alone implies a quantum
advantage, but it gives a reproducible pre-submission filter for deciding which
integrals are structurally appropriate for amplified QAE and which should remain
in the direct Grover--Rudolph sampling regime.

The Phase 12 local QAE-amplitude audit adds a second filter: the scalar
amplitude \(a=I_p(f)\) and the selected amplification schedule must be
informative. The low-angle upper-half tests for the \(\sin^2(\pi x)\) profiles
have \(a=1/2\), so \(p_k(a)=1/2\) for every Grover level and they are not useful
QAE estimation tests despite their favorable angle class. The first meaningful
non-degenerate augmented-QAE candidate is instead the eight-cell
\(\sin^2(\pi x)\) law with \(f(x)=\sin(\pi x)\). This pair remains simulator
only for now; a backend-aware noisy audit should precede any QPU submission.

That backend-aware audit has now been performed in Phase 13. On both
\texttt{ibm\_kingston} and \texttt{ibm\_fez}, the noisy model drives the
augmented-QAE marked probabilities toward approximately \(1/2\), even though
the ideal values for \(k=0\) and \(k=1\) are well separated. The full
\(K=\{0,1,2\}\) schedule is also not viable under the current cost thresholds.
Thus the dual angle-structure workflow has served its intended purpose: it
identified a mathematically meaningful QAE candidate, then rejected it before
QPU spending once backend-aware noise was included. For the next hardware
campaign, the evidence still favors direct Grover--Rudolph distribution
sampling and quantity-of-interest postprocessing over amplified QAE.

Phase 14 begins that next direction by testing direct integration scaling
locally from \(n=4\) to \(n=6\). In a local \texttt{cz}-basis audit, all three
families reach 64 midpoint cells without exceeding the provisional cost
thresholds; the \(n=6\) cases require 62 two-qubit gates in the best local
transpilation. This makes \(n=6\) a plausible target for backend-aware auditing,
but not yet for QPU execution. The Phase 9 and Phase 10 evidence still matters:
the affine family is the most credible scaling candidate, while the localized
beta-like family should be treated primarily as a stress diagnostic unless a
backend noise audit shows a clear improvement.

The backend-aware Phase 14 audit refines this decision. On
\texttt{ibm\_kingston}, all three \(n=6\) candidates pass the pre-submission
thresholds, while on \texttt{ibm\_fez} only the localized beta-like profile
passes. The recommended next QPU microcampaign is nevertheless the affine
64-cell law on \texttt{ibm\_kingston}, because it is the direct continuation of
the 32-cell affine result that already passed the Phase 9--10 paper-quality
filters. The beta-like 64-cell law should be reserved for a later stress test,
not used as the main scaling claim.

The executed 64-cell affine campaign confirms this nuanced view. At the
distribution level the TVD rises to 0.058478, close to the Phase 14 acceptance
boundary, but the maximum cell deviation remains small. At the integration
level the smooth quantity \(I_p(\sin(\pi x))\) improves to absolute error
0.000752 and passes interval coverage, making it the strongest current
paper-quality direct-integration result. In contrast, the upper-half mass, the
mean, the second moment, and the call payoff show clear residual bias. Future
claims should therefore be quantity-specific: the experiment supports direct
Grover--Rudolph integration for selected smooth quantities on 64 cells, not a
uniform backend-independent integration advantage.

The Phase 15 Gaussian test reinforces the same conclusion from the opposite
direction. A smoother centered law with a much shallower circuit improves the
upper-half mass and the mean, but it worsens \(I_p(\sin(\pi x))\) and does not
improve the global distribution TVD. The finite-law experiments should
therefore be presented as a controlled hardware study of probability-law
preparation and quantity-specific integration, not as a single monotone scaling
story.

This also clarifies the transition to the next phase. The direct
Grover--Rudolph workflow estimates \(I_p(f)\) for a prepared law \(p\). The
angle-structure paper studies a different numerical-integration object:
quadrature-rule QAE for \(I[g]=\int_0^1 g(x)\,dx\), with explicit left,
midpoint, right, and Simpson rules. The next experimental block should use the
finite-law results as hardware context, but should return to the paper's
quadrature formulation in order to test the encoding-complexity theory
directly.

That transition has now started with the Phase 15b midpoint-quadrature QAE
microcampaign on \texttt{ibm\_kingston}. The restored Workflow-A scripts
separate \(Q_n[g]\), \(I[g]\), MLAE estimation error, total error, circuit
depth, two-qubit gates, shots, and oracle-query count. The hardware results are
encouraging: \(g_1\) gives total error 0.001172, \(g_2\) gives total error
0.000140, and the smooth low-amplitude bump gives total error 0.001568. The
shifted affine stress case is less accurate, with error 0.010640, and should
remain a watch case. The main conclusion is therefore positive but precise:
quadrature-rule MLAE on current IBM hardware is viable for the low-degree
calibration hierarchy and at least one smooth low-amplitude stress integral,
but backend-aware admission remains necessary before extending to larger
grids, Simpson combinations, or non-special affine amplitudes.

The \texttt{ibm\_aachen} repeat sharpens this point. Although
\texttt{ibm\_aachen} had better average two-qubit and readout calibration
figures than \texttt{ibm\_kingston} at submission time, it did not improve the
microcampaign globally. Kingston remains better for \(g_1\), \(g_2\), and the
shifted affine stress case; Aachen improves only the smooth bump slightly. The
large affine failure on Aachen, driven by the \(k=1\) amplified probability,
shows that backend choice for QAE cannot be made from aggregate calibration
numbers alone. The practical rule for the next phase is to run a backend-aware
noisy/layout audit for each proposed amplified circuit and to treat
cross-backend repeats as validation, not as a substitute for circuit-level
admission tests.
